% All-rank structural cancellation; not a solution of arbitrary shallow deformations.
\begin{proposition}[Proved: trace cancellation for all commuting shallow cubic perturbations]
\label{C:shallow_commuting_trace}
Let $A=\mathbb Z_3[\zeta]$, $\pi=\zeta-1$, and
$k=A\ltimes J$, where $J$ is any $\mathbb F_3$-vector space.
Let $L$ be a finite free $A[C_3]$-lattice with generator $T_0$.
Suppose $U\in\operatorname{End}_A(L)\otimes_AJ$ commutes
with the residual action $\overline T_0$; equivalently,
$T_0U=UT_0$ as maps into $L\otimes_AJ$.
Then $T_0+U$ defines a $C_3$-action on $L\otimes_Ak$, and
\[
 [(L\otimes_Ak,T_0+U)]-[(L\otimes_Ak,T_0)]
       =\delta(\operatorname{tr}U),\qquad
 \delta(x)=[k_{1+x}]-[k].                              \tag{C.1.34}
\]
Its induction to $S_3$ is zero.
The same conclusion holds if an infinitesimal change of basis
makes the perturbation commute with the residual action.
This covers arbitrary ranks and all constant integral lattice
types, but does not include arbitrary noncommuting perturbations.
\end{proposition}
\begin{proof}
The products of two $J$-valued matrices vanish. Commutation
and $3J=0$ therefore give $(T_0+U)^3=1$.
For a fixed lattice write its relative class as $d_L(U)$,
and put $z=[k_\zeta]$.
For any two perturbations $U,V$ each commuting with the residual action, the matrices
\[
 E_U=\begin{pmatrix}T_0+U&I\\0&\zeta T_0\end{pmatrix},
 \qquad P=\begin{pmatrix}I&0\\V-U&I\end{pmatrix}
\]
satisfy
\[
 E_U^3=I,\qquad
 P^{-1}E_UP=
 \begin{pmatrix}T_0+V&I\\0&\zeta T_0+U-V\end{pmatrix}.
                                                               \tag{C.1.35}
\]
Indeed the constant off-diagonal cubic sum is zero, while
its linear part is $[\overline T_0,U]=0$; the lower block
in the conjugacy vanishes by the same commutation and
$(\zeta-1)J=J^2=0$.
The two free-block exact sequences give
\[
 d_L(U)=d_L(V)+z\,d_L(U-V).
\]
Putting $U=0$ and applying the relation twice gives
$z^2d_L(V)=d_L(V)$. Since $z^3=1$, it follows that
$zd_L(V)=d_L(V)$. Thus $d_L$ is additive on the space of
commuting perturbations, and is $\mathbb F_3$-linear.

Decompose $L$ into the nine types of
Proposition~\ref{C:cyclotomic_C3_lattices}.
Each block of a commuting $U$ is an intertwiner of the
corresponding residual actions. Additivity permits treating
one block at a time. An off-diagonal block gives a triangular
action with constant diagonal blocks and hence relative
class zero. It remains to consider a single indecomposable
constant lattice.

For a character line the assertion is precisely the rank-one
formula and $z\delta(x)=\delta(x)$ from
Proposition~\ref{C:cyclotomic_torsion_cancellation}.
For a pair gluing its residual action is $I+N$, $N=E_{12}$.
Its commutant consists of $aI+bN$. The $bN$ term is an
infinitesimal conjugation, since $N=[N,E_{22}]$.
The scalar term preserves an integral character flag and
therefore has relative class $2\delta(a)=\delta(2a)$.
For the regular lattice every deformation is isomorphic to
the regular module: lift a regular generator, and the resulting
map from $k[C_3]$ is an isomorphism modulo $J$, hence an
isomorphism. The residual commutant is
$\mathbb F_3[N]/(N^3)$, whose matrices have trace zero.

For either remaining rank-three type $L_{012}$ or $Q$,
choose a residual basis with $N=E_{12}$ and third vector $Y$.
Every commuting perturbation has the form
\[
 U=\begin{pmatrix}a&b&c\\0&a&0\\0&d&e\end{pmatrix},
 \qquad a,b,c,d,e\in J.
\]
The three off-diagonal terms are infinitesimal conjugations:
$E_{12},E_{13},E_{32}$ are commutators with $N$.
The remaining terms are $aI+(e-a)E_{33}$.
Every scalar perturbation on a rank-three lattice has class
$3\delta(a)=0$: a saturated integral character flag exists
over the DVR and is preserved by a scalar.

For $L_{012}$, the primitive invariant vector $\pi e_0$ has
residue in $\ker N$ outside $\operatorname{im}N$. Take it
as $Y$, and lift the other two residual chain vectors to an
$A$-basis. The perturbation $xE_{33}$ preserves this free
character line, acts on it by $1+x$, and is zero on the
free quotient. Hence $d_{L_{012}}(xE_{33})=\delta(x)$.

For $Q$, the coordinate map $(a_0,a_1,a_2)\mapsto a_0$
is a surjective equivariant map to the trivial character.
Its reduction is nonzero on $\ker N/\operatorname{im}N$.
Choose $Y$ mapping to one and complete it by a residual
Jordan chain in its kernel. With an $A$-basis lifting this
choice, the perturbation $xE_{33}$ is zero on that free
kernel and acts by $1+x$ on the quotient. This again gives
$d_Q(xE_{33})=\delta(x)$.
For an explicit check, use the basis
$(e_0-e_2,e_1-e_2,\pi e_2)$ of $Q$: its residual difference
has last row $(1,2,0)$, and the first-coordinate functional
has row $(1,0,0)$, which is nonzero on the residual kernel
line represented by $(1,1,0)$.

Thus in both cases the relative class is $\delta(e-a)$.
This is $\delta(\operatorname{tr}U)$, since
$2a+e=e-a$ in $J$. Together with the preceding types and
block additivity this proves (C.1.34) in every rank.
Induction kills its right side by
Proposition~\ref{C:cyclotomic_torsion_cancellation}.
Finally an infinitesimal conjugation changes $U$ by a
commutator with $\overline T_0$, whose trace is zero, and
does not change the representation class. This proves the
last assertion.
\end{proof}
